\documentclass[10pt,a4paper]{amsart}
\usepackage[margin=19mm]{geometry}
\usepackage{amsmath,amssymb,mathtools}
\usepackage[scr=rsfs,cal=euler]{mathalfa}
\usepackage{xfrac}
\usepackage[unicode,hidelinks,bookmarks=false]{hyperref}
\newcommand{\ie}{i.e.\ }
\newcommand{\cf}{cf.\ }
\newcommand{\resp}{respectively\ }
\newcommand{\Subsection}{\subsection}
\providecommand{\nobreakdash}{\nobreak\hbox{-}}
\setlength{\emergencystretch}{2em}
\multlinegap0pt
\usepackage{amssymb}
\usepackage{mathtools}
\usepackage{enumerate}
\newtheorem{lemma}{Lemma}
\title{A priori estimates for solutions of degenerate fully nonlinear elliptic equations with $L^p$ data}
\author{Hongsoo Kim, Se-Chan Lee}
\date{Source: arXiv:2605.20650v1 (CC BY 4.0)}
\begin{document}
\maketitle

\noindent\emph{Source and licence.} A priori estimates for solutions of degenerate fully nonlinear elliptic equations with $L^p$ data. Version arXiv:2605.20650v1 (CC BY 4.0), distributed by arXiv under the Creative Commons Attribution 4.0 International licence, http://creativecommons.org/licenses/by/4.0/, as recorded in the arXiv licence metadata for this exact version and shown on the abstract page https://arxiv.org/abs/2605.20650v1. Full licence notice: \texttt{licenses/arxiv-2605.20650-CC-BY-4.0.txt}.\par\smallskip

\noindent\textit{Editorial scope.} Counted unit: the article's lemma A01 (source0.tex lines 435--448). The article's complete original proof of it is delivered with it (source0.tex lines 450--510, one \texttt{proof} environment of 61 lines). No mathematics is written, repaired or re-derived: the statement and the proof are the article's own lines, in the article's own notation, and no retained line is substituted or re-flowed.  No further source-local context is needed: every cross-reference of every retained line resolves inside the two delivered spans.  Nothing was omitted from any retained span, so the package carries no omission record.  No retained line contains a citation, so the wrapper carries no bibliography and no bibliography entry.  No retained line contains a figure environment, an image-inclusion command or drawing code, so no image is delivered and the package declares no asset dependency.  Editorial formatting only, in addition: an AMS \texttt{amsart} preamble that restates the article's own declarations and macros used by the retained lines, namely the environments \texttt{lemma} on separate counters, together with the standard packages those lines need.  The retained environments print in this document as Lemma 1 in order of appearance, whereas the article prints the same items with the numbering of its own full document.  The complete native source of the article as distributed in the arXiv e-print for this exact version is bundled unchanged as \texttt{sources/201006/source0.tex}.
\begin{lemma}\label{A01}
    There exist small $\delta_0>0$ and large $K>1$, $A_0>1$ such that if  $u \in C(B_3)$ satisfies
    \begin{equation*}
    \left\{
    \begin{aligned}
        u &\geq 0 &&\text{in } B_3,\\
        \mathcal{M}^-(D^2u)-\Lambda|Du| &\leq f &&\text{in } \{|Du| \leq A_0\}\cap B_3, \\
        \|f\|_{L^n(B_3)} &\leq \delta_0, \\
        u&>K &&\text{in } B_{1/4},
    \end{aligned}
    \right.
    \end{equation*}
    then $u>1$ in $B_1$.
\end{lemma} 

\begin{proof}
By standard approximation again, we may assume that $u\in C^2(B_3)\cap C(\overline{B_3})$. We use the contradiction by assuming that $u(x_0) \leq 1$ for some $x_0 \in B_{1}$.
We define the following barrier function
\begin{equation*}
    \Phi(z) =  \frac{M}{p}(|z|^{-p}-2^{-p}) \quad \text { in } \mathbb{R}^n\setminus B_{1/8},
\end{equation*}
for some $p>1$ and $M>1$ to be determined later, and extend $\Phi (>0)$ smoothly inside $B_{1/8}$. For $z \in B_4\setminus B_{1/8}$, a direct calculation shows
\begin{equation} 
\begin{aligned}\label{barrierPDE}
    \mathcal{M}^-(D^2\Phi) -\Lambda|D\Phi| &= M |z|^{-p-2}(\lambda(p+1)- \Lambda(n-1) - \Lambda|z|) \\
    &\geq 1, 
\end{aligned}    
\end{equation}
provided that $p=p(n,\lambda,\Lambda)>1$ and $M=M(n,\lambda,\Lambda)>1$ are sufficiently large.

We choose $M$ further larger to ensure that $\Phi >2$ in $B_{3/2}$.
Consider the vertex set $V = B_{1/8}$ and we slide $\Phi(\cdot-y)$ with vertex $y \in V$ from below, until it touches the graph of $u$ for the first time at a point $x \in \overline{B_3}$. Let $T$ denote the set of corresponding contact points.

For $y \in V$, we observe that
\begin{equation*}
    u(x_0)-\Phi(x_0-y) < 1-2=-1, \quad \text{since $x_0 \in B_1$}.
\end{equation*}
If $x\in \partial B_3$, then since $u \geq 0$  and $\partial B_{3} \subset \{z:\Phi(z-y) \leq 0\}$, we find that
\begin{equation*}
    u(x) - \Phi(x-y) \geq 0.
\end{equation*}
Therefore, we conclude that the minimum of $u-\Phi(\cdot-y)$ cannot be attained on the boundary $\partial B_3$ and $T \subset \{ u<K\} \cap B_3$ for $K=\|\Phi\|_{L^\infty(\mathbb{R}^n)}$.
Recalling that $u>K$ in $B_{1/4}$, we have $T \subset \{ u<K\} \cap (B_3 \setminus B_{1/4})$.
Moreover, we find $1/8<|x-y| < 4$ for $x \in T$ and $y \in V$.
Note that for any $x \in T$, we get
\begin{align} \label{Ducond2}
    Du(x) &= D\Phi(x-y),\\ \label{D2ucond2}
    D^2u(x) &\geq D^2\Phi(x-y).
\end{align}
Choosing $A_0=\sup_{B_4\setminus B_{1/8}}|D\Phi|$, we obtain $T \subset \{|Du| \leq A_0\}$, and so the inequality
\begin{equation*}
    \mathcal{M}^-(D^2u)-\Lambda|Du| \leq f
\end{equation*}
is satisfied in $T$. Therefore, for $x \in T$, we have
\begin{equation} \label{uequ}
    1 \leq \mathcal{M}^-(D^2\Phi) -\Lambda|D\Phi| \leq \mathcal{M}^-(D^2u) -\Lambda|Du| \leq f(x).
\end{equation}
By combining \eqref{D2ucond2} and \eqref{uequ}, we find
\begin{equation*}
    |D^2u(x)| \leq C(|D^2\Phi(x-y)| + |Du(x)| +f(x)) \leq Cf(x),
\end{equation*}
where we used $|D^2\Phi(x-y)| \leq C=C(n, \lambda, \Lambda)$.
By differentiating \eqref{Ducond2} with respect to $x$ together with the fact that $D^2\Phi(\cdot)$ is (uniformly) invertible in $B_4 \setminus B_{1/8}$, we get
\begin{equation*}
	D_xy= I- (D^2 \Phi)^{-1}(x-y)  \cdot D^2u(x),
\end{equation*}
which implies that
\begin{equation*}
   |\det D_xy| \leq C|f(x)|^n.
\end{equation*}
Applying the area formula, we arrive at
\begin{equation*}
    |B_{1/8}| = |V| =\int_{T}|\det D_xy|\,dx \leq C\|f\|_{L^n(T)}^n \leq C\delta_0^n.
\end{equation*}
Choosing $\delta_0$ sufficiently small yields a contradiction.
\end{proof}

\end{document}
